\documentclass[12pt]{article}
\usepackage[utf8]{inputenc}
\usepackage{amsmath, amssymb, amsfonts}
\usepackage{geometry}
\usepackage{booktabs}
\usepackage{enumitem}
\usepackage{siunitx}
\geometry{margin=1in}
\setlist[enumerate]{leftmargin=*}
\newcommand{\bn}{\bigskip\noindent}
\newcommand{\mc}[1]{\mathcal{#1}}

\title{Solutions -- Practice Problem Set 2\\ \vspace{15pt} \large ECON 101 -- Macroeconomics}
\author{}
\date{Fall 2026}

\begin{document}
\maketitle



\begin{enumerate}

\item
A study that makes coffee more attractive increases consumers' willingness to pay at each price.\\
\textbf{Graph:} Demand shifts right from $D_0$ to $D_1$; supply $S$ is unchanged.\\
\textbf{New equilibrium.} The intersection of $D_1$ and $S$ implies a higher equilibrium price $P_1>P_0$ and higher equilibrium quantity $Q_1>Q_0$.\\
\textbf{Explanation:} At the original price $P_0$, the rightward shift creates excess demand, bidding up price. The price rises until the new intersection clears the market.

\item
 A cost-reducing technology raises firms' willingness to supply at each price.\\
\textbf{Graph:} Supply shifts right from $S_0$ to $S_1$; demand $D$ is unchanged.\\
\textbf{New equilibrium:} The intersection of $S_1$ and $D$ implies a lower equilibrium price $P_1<P_0$ and higher equilibrium quantity $Q_1>Q_0$.\\
\textbf{Explanation:} At the original price $P_0$, the rightward shift creates excess supply, pushing price down. The price falls until the new intersection clears the market.

\item
\[
Q_d=100-2P, \qquad Q_s=20+3P.
\]
\textbf{Equilibrium.} Set $Q_d=Q_s$:
\[
100-2P=20+3P \;\Rightarrow\; 80=5P \;\Rightarrow\; P^\ast=16.
\]
\[
Q^\ast=100-2(16)=100-32=68 \quad \text{(checks: }20+3\cdot 16=20+48=68\text{)}.
\]
\textbf{At $P=10$.} 
\[
Q_d=100-2(10)=80, \qquad Q_s=20+3(10)=50.
\]
\textbf{Result.} \;Shortage of $80-50=30$ units.

\item
\[
Q_d=100-2(20)=100-40=60, \qquad Q_s=20+3(20)=20+60=80.
\]
\textbf{Surplus.} \; $80-60=20$ units.\\
\textbf{Binding?} Yes. The floor $20$ exceeds the equilibrium price $16$, so it prevents the market from reaching $P^\ast$ and creates excess supply.

\item
\[
Q_d=80-2P,\qquad Q_s=10+2P.
\]
\textbf{Equilibrium without ceiling.} Set $Q_d=Q_s$:
\[
80-2P=10+2P \;\Rightarrow\; 70=4P \;\Rightarrow\; P^\ast=17.5.
\]
\[
Q^\ast=80-2(17.5)=80-35=45 \quad \text{(checks: }10+2(17.5)=10+35=45\text{)}.
\]
\textbf{With ceiling } $P^c=15$:
\[
Q_d=80-2(15)=80-30=50,\qquad Q_s=10+2(15)=10+30=40.
\]
\textbf{Shortage.} \; $50-40=10$ units.\\
\textbf{Possible unintended consequence.} Reduced maintenance and quality of rental units as landlords face lower effective prices and weaker incentives to invest.












\item \textbf{Expenditure Approach}

Data (billions): $C=1400$, $I=500$, $G=600$, $X=400$, $M=450$.

\bn \textit{(a) GDP using expenditure approach}
\[
\text{GDP} = C + I + G + (X - M)
= 1400 + 500 + 600 + (400 - 450)
= 1400 + 500 + 600 - 50
= \boxed{2450}
\]
(billions of dollars)

\bn \textit{(b) NNP and GNP given depreciation $=200$, net factor payments from abroad $=\text{NFP}=-50$}
\[
\text{GNP} = \text{GDP} + \text{NFP} = 2450 + (-50) = \boxed{2400}
\]
\[
\text{NNP} = \text{GNP} - \text{Depreciation} = 2400 - 200 = \boxed{2200}
\]
(all in billions of dollars)

\bigskip

\item \textbf{Income Approach}

Data (billions): Wages \& Salaries $=2200$, Rental Income $=300$, Corporate Profits $=400$, Net Interest $=200$, Indirect Taxes less Subsidies $=250$, Depreciation $=150$.

\bn \textit{(a) GDP via income approach}
\[
\text{GDP} =
\underbrace{(2200 + 300 + 400 + 200)}_{\text{factor incomes}}
+ \underbrace{250}_{\text{indirect taxes} - \text{subsidies}}
+ \underbrace{150}_{\text{depreciation}}
= 3500
\]
\[
\boxed{\text{GDP} = 3500 \text{ (billions)}}
\]

\bn \textit{(b) Compare to expenditure GDP of \$3{,}350b}
\[
\text{Statistical discrepancy} = 3500 - 3350 = \boxed{150 \text{ (billions)}}
\]
Interpretation: measurement and timing differences across sources leave a positive discrepancy of \$150b.

\bigskip

\item \textbf{Real vs Nominal GDP}

Two goods: bread and smartphones.

\begin{center}
\begin{tabular}{lcccc}
\toprule
Year & Bread (Qty) & Bread Price & Smartphones (Qty) & Phone Price \\
\midrule
2023 & 100 & \$2.00 & 50 & \$400 \\
2024 & 110 & \$2.20 & 60 & \$380 \\
\bottomrule
\end{tabular}
\end{center}

\bn \textit{(a) Nominal GDP}
\[
\text{Nominal GDP}_{2023} = 100\times 2 + 50\times 400 = 200 + 20000 = \boxed{20200}
\]
\[
\text{Nominal GDP}_{2024} = 110\times 2.20 + 60\times 380 = 242 + 22800 = \boxed{23042}
\]

\bn \textit{(b) Real GDP (base year = 2023)}
\[
\text{Real GDP}_{2023}^{(2023\text{ prices})} = \boxed{20200}
\]
\[
\text{Real GDP}_{2024}^{(2023\text{ prices})}
= 110\times 2.00 + 60\times 400 = 220 + 24000 = \boxed{24220}
\]

\bn \textit{(c) GDP deflator, 2024}
\[
\text{Deflator}_{2024} = \frac{\text{Nominal}_{2024}}{\text{Real}_{2024}} \times 100
= \frac{23042}{24220}\times 100 \approx \boxed{95.14}
\]
(2023 base $=100$ implies an overall price decline driven by cheaper smartphones.)

\end{enumerate}



\begin{enumerate}[resume]
\item \textbf{CPI Calculation}

Basket: 10 bread, 5 movie tickets.

\bn \textit{(a) CPI in 2023 and 2024}
\[
\text{Cost}_{2023} = 10\times 2 + 5\times 10 = 20 + 50 = 70
\quad\Rightarrow\quad \text{CPI}_{2023} = \boxed{100}
\]
\[
\text{Cost}_{2024} = 10\times 2.20 + 5\times 12 = 22 + 60 = 82
\quad\Rightarrow\quad \text{CPI}_{2024} = \frac{82}{70}\times 100 \approx \boxed{117.14}
\]

\bn \textit{(b) Inflation rate (CPI), 2023 to 2024}
\[
\pi = \frac{\text{CPI}_{2024} - \text{CPI}_{2023}}{\text{CPI}_{2023}}\times 100
= \frac{117.14 - 100}{100}\times 100 \approx \boxed{17.14\%}
\]

\bn \textit{(c) Compare CPI inflation with GDP deflator from Q3}
\[
\text{GDP deflator change} \approx 95.14 - 100 = -4.86 \Rightarrow \text{deflation of } \boxed{4.86\%}
\]
Differences arise because:
\begin{itemize}
    \item \textbf{Coverage}: CPI tracks consumer purchases, includes imports; GDP deflator covers domestically produced final goods and services, excludes imports and includes exports.
    \item \textbf{Weights}: CPI fixes a basket; GDP deflator uses current production weights.
    \item \textbf{Relative price movements}: Large smartphone price declines lower the deflator more, while the CPI basket assigns different weights and may reflect higher movie-ticket inflation.
\end{itemize}

\bigskip

\item \textbf{Nominal vs Real Wages}

CPI: $100 \to 110$. Nominal wage: \$20 (2023) $\to$ \$21 (2024).

\bn \textit{(a) Real wage in 2024 (base 2023=100)}
\[
w^{\text{real}}_{2024} = \frac{w^{\text{nom}}_{2024}}{\text{CPI}_{2024}/100}
= \frac{21}{1.10} \approx \boxed{\$19.09}
\]
(with base-year purchasing power)

\bn \textit{(b) Did purchasing power rise?}
In 2023, real wage was \$20 (base). In 2024, it is about \$19.09, so purchasing power \emph{fell}. The nominal raise did not keep up with inflation.

\end{enumerate}

\end{document}
